\documentclass[tikz,border=5pt]{standalone}
\usepackage{ctex}
\usepackage{amsmath,amssymb}
\usetikzlibrary{arrows.meta,positioning,fit}
% Shared visual language for LFS architecture diagrams.
\RequirePackage{../../mymath}

\definecolor{LFSBlue}{HTML}{2563EB}
\definecolor{LFSBlueSoft}{HTML}{DBEAFE}
\definecolor{LFSGreen}{HTML}{15803D}
\definecolor{LFSGreenSoft}{HTML}{DCFCE7}
\definecolor{LFSOrange}{HTML}{C2410C}
\definecolor{LFSOrangeSoft}{HTML}{FFEDD5}
\definecolor{LFSRed}{HTML}{B91C1C}
\definecolor{LFSRedSoft}{HTML}{FEE2E2}
\definecolor{LFSPurple}{HTML}{7E22CE}
\definecolor{LFSPurpleSoft}{HTML}{F3E8FF}
\definecolor{LFSGray}{HTML}{475569}
\definecolor{LFSGraySoft}{HTML}{F1F5F9}

\tikzset{
  lfs picture/.style={
    font=\small,
    node distance=8mm and 12mm,
    >=Latex,
    every path/.style={line cap=round, line join=round}
  },
  block/.style={
    draw=LFSBlue,
    fill=LFSBlueSoft,
    rounded corners=2pt,
    line width=0.8pt,
    align=center,
    inner xsep=7pt,
    inner ysep=5pt,
    minimum height=10mm
  },
  compute/.style={
    block,
    draw=LFSPurple,
    fill=LFSPurpleSoft
  },
  storage/.style={
    block,
    draw=LFSGreen,
    fill=LFSGreenSoft
  },
  interface/.style={
    block,
    draw=LFSOrange,
    fill=LFSOrangeSoft
  },
  risk/.style={
    block,
    draw=LFSRed,
    fill=LFSRedSoft
  },
  neutral/.style={
    block,
    draw=LFSGray,
    fill=LFSGraySoft
  },
  decision/.style={
    diamond,
    draw=LFSOrange,
    fill=LFSOrangeSoft,
    line width=0.8pt,
    align=center,
    inner sep=2pt,
    aspect=2
  },
  group/.style={
    draw=LFSGray,
    rounded corners=3pt,
    dashed,
    line width=0.7pt,
    inner sep=6pt
  },
  arrow/.style={
    -{Latex[length=2.4mm,width=1.6mm]},
    draw=LFSGray,
    line width=0.9pt
  },
  data arrow/.style={
    arrow,
    draw=LFSBlue
  },
  control arrow/.style={
    arrow,
    draw=LFSOrange,
    dashed
  },
  residual/.style={
    arrow,
    draw=LFSGreen,
    rounded corners=4pt
  },
  label text/.style={
    font=\scriptsize,
    text=LFSGray,
    fill=white,
    inner sep=1.5pt
  },
  title/.style={
    font=\bfseries,
    text=LFSGray,
    align=center
  }
}


\begin{document}
\begin{tikzpicture}[my picture, >=Latex, node distance=8mm]

\tikzset{
  biasbox/.style={draw=MyRed, fill=MyRedSoft, rounded corners=3pt, font=\small\bfseries, align=center, inner sep=5pt, minimum width=46mm, minimum height=14mm},
  fixbox/.style={draw=MyGreen, fill=MyGreenSoft, rounded corners=3pt, font=\small, align=center, inner sep=5pt, minimum width=46mm, minimum height=14mm},
  title/.style={font=\normalsize\bfseries, text=MyGray, align=center}
}

\node[title] (t1) at (2.3, 4.5) {大模型裁判系统性偏置 (Systematic Biases)};
\node[title] (t2) at (7.8, 4.5) {受控校准与缓解协议 (Mitigations)};

% Bias 1: Position Bias
\node[biasbox] (b1) at (2.3, 3.2) {
  \textbf{位置偏置 (Position Bias)}\\
  \scriptsize 倾向于给排在前面的候选打更高分
};
\node[fixbox] (f1) at (7.8, 3.2) {
  \textbf{双向位置对调 (Position Swap)}\\
  \scriptsize 交换候选次序盲评，仅双向一致才采纳
};
\draw[->, draw=MyOrange, line width=1pt] (b1) -- (f1);

% Bias 2: Verbosity Bias
\node[biasbox] (b2) at (2.3, 1.6) {
  \textbf{冗长偏置 (Verbosity Bias)}\\
  \scriptsize 误将篇幅长、格式花哨当做高回答质量
};
\node[fixbox] (f2) at (7.8, 1.6) {
  \textbf{长度归一化与约束量表}\\
  \scriptsize 强制按事实论断原子切分，惩罚无意义冗余
};
\draw[->, draw=MyOrange, line width=1pt] (b2) -- (f2);

% Bias 3: Self-enhancement Bias
\node[biasbox] (b3) at (2.3, 0.0) {
  \textbf{自我偏好 (Self-Enhancement)}\\
  \scriptsize 显著偏好自身或同一模型家族生成的文风
};
\node[fixbox] (f3) at (7.8, 0.0) {
  \textbf{多模型裁判委员会 (Panel of LLMs)}\\
  \scriptsize 采用异构架构模型交叉打分与加权仲裁
};
\draw[->, draw=MyOrange, line width=1pt] (b3) -- (f3);

% Bias 4: Prompt Injection
\node[biasbox] (b4) at (2.3, -1.6) {
  \textbf{裁判提示注入 (Judge Injection)}\\
  \scriptsize 待评测文本包含对抗指令操控裁判输出
};
\node[fixbox] (f4) at (7.8, -1.6) {
  \textbf{输入结构隔离与沙箱化}\\
  \scriptsize 候选文本放入独立变量块，格式强校验
};
\draw[->, draw=MyOrange, line width=1pt] (b4) -- (f4);

\end{tikzpicture}
\end{document}
